\section{Z=1 empirical triangulation of \texorpdfstring{$r_e/r_0$}{r\_e/r\_0}}
\label{sec:z1-triangulation}

\paragraph{Result.}
\[
\boxed{\;
  \frac{r_e}{r_0} \;=\; 0.4994205099128317 \;\pm\; 2.50\times 10^{-13}
\;}\qquad \text{(PR \#62, joint best-fit across six observables).}
\]

\paragraph{Method (full derivation: \S\ref{sec:eq-triangulation}).}
Per \texttt{10\_CrossComparison.md}~\S2, the six $g_s$-dependent observables
each enter as $\mathrm{anchor}_i \cdot (g_r(x)/g_s^{\rm meas})^{n_i}$ with
$n_i \in \{1, 2\}$. Two weighting passes are reported:
\begin{itemize}[topsep=2pt,itemsep=0pt]
  \item \textbf{Pass~A} (measurement-$\sigma$ only): $r_{\rm opt} =
        0.4994061257148855$ at $\chi^2 \approx 2.97\times 10^{16}$.
        Unphysical: pulls $g_e(r_{\rm opt})$ off measured by
        $5.76\times 10^{-5}$ to compensate the hyperfine residual that the
        leading-$(g_s/{-2})^n$ model cannot fit by any choice of $r$ alone.
  \item \textbf{Pass~B} (measurement-$\sigma$ + framework-floor noise):
        $r_{\rm opt} = \mathbf{0.4994205099128317}$ at
        $\chi^2 \approx 3.99958 \approx 4$ (the number of framework-floor
        residuals at $\sim$one-floor). Matches branch~(c) to 16 sig figs.
\end{itemize}
Pass~B is the honest weighting. Stretched-fit flag: \textbf{none}.

\paragraph{Closed-form algebraic inversion (the \S III.D-extension).}
\texttt{Roadmapping/Equation\_Verification/Dual\_Relativistic\_Quantum\_Mechanics\_I.md}
\S III.D-extension establishes the closed-form
\[
\boxed{\;
  \frac{r_e}{r_0} \;=\; \frac{2 - a_e}{2(2 + a_e)}
\;}
\]
obtained by inverting $g_r(r_e/r_0) = -2(1 + a_e)$. With CODATA-full
$a_e^{\rm expt}$, this gives $r_e/r_0 = 0.499\,420\,509\,913\,18$, matching
the triangulated optimum to $3.45\times 10^{-13}$ (within precision floor
$\sigma_r = 2.5\times 10^{-13}$).

\paragraph{Honest framing (verbatim from DRQM~I~\S III.D-extension).}
This is \emph{reproduction-by-inheritance under hypothesis~(ii)}: QED supplies
$a_e(\alpha)$; the framework supplies the (III.22) cutoff--anomaly
identification and the algebraic inversion. It is \textbf{not} an independent
dual-framework derivation of $a_e$. The \S III.D Schwinger identification
subsection (PR~\#70) characterises the same algebraic operation as the
back-fit identity from the Candidate-3 perspective; both subsections together
give the full picture.

\paragraph{Finding 2 verdict trajectory (verbatim from
\texttt{FINDINGS\_for\_author\_review.md}:209).}
\[
\mathbin{\hbox{\color{red}\(\bullet\)}}\,\text{(was)}
\;\longrightarrow\;
\mathbin{\hbox{\color{yellow!90!black}\(\bullet\)}}\,\text{characterised (\#61)}
\;\longrightarrow\;
\mathbin{\hbox{\color{yellow!90!black}\(\bullet\)}}/
\mathbin{\hbox{\color{green!60!black}\(\bullet\)}}\,
\text{at hypothesis-(ii) framework precision (\#64)}
\]
\textbf{Unconditional} $\mathbin{\hbox{\color{green!60!black}\(\bullet\)}}$
requires hypothesis-(i) re-derivation (proper-time photon propagator +
bound-state propagator in the (II.3) Pauli kernel + mass-renormalisation
prescription at the cutoff), tracked in
\textbf{\href{https://github.com/temoTxt/PyPhysics/issues/75}{issue \#75}}.

\paragraph{Hypothesis split (the load-bearing question).}
\begin{enumerate}[topsep=2pt,itemsep=0pt,label=(\roman*)]
  \item the dual-framework proper-time one-loop vertex (with the dual photon
        propagator and bound-state structure) is computed \emph{independently}
        and shown to reproduce $a_e^{(1)} = \alpha/(2\pi)$;
  \item the dual-framework one-loop vertex correction \emph{reproduces} the
        standard-QED Schwinger anomaly identically by inheritance, with all
        dual structure absorbed into the matter-side (II.3) kernel which at
        one-loop is formulation-independent.
\end{enumerate}
The closed-form inversion + CODATA-$a_e$ chain delivers the triangulated value
under either reading; only (i) constitutes a derivation \emph{within} the
framework.

\paragraph{Variational / operator-coefficient route (Candidate~2).}
Route~Z's framework-internal closure (closure 7c on the spin-magnetic operator
coefficient, tree-level target $g = -2$) yields
$r_e/r_0 = 1/2$ exactly, matching FoundationsII-Classical \S2.2's
critical-point identification. The discrepancy with the triangulated value
$0.4994205099128317$ is $5.79\times 10^{-4}$ --- exactly the size of the
Schwinger one-loop QED correction $\alpha/(2\pi) \approx 1.16\times 10^{-3}/2
\approx 5.8\times 10^{-4}$. The Schwinger 1-loop refinement reproduces the
triangulated value to $\sim 10^{-7}$ (Karplus--Kroll-level), but requires QED
input \emph{external} to the framework. Full derivation:
\S\ref{sec:eq-variational} and \S\ref{sec:eq-schwinger-residual}.

\authorreview{(1)~Does the Pass~B weighting match your intent? (2)~Is
``reproduction-by-inheritance under hypothesis~(ii)'' the framework's
intended position for the closed-form result, or is hypothesis~(i)
(issue~\#75) required? (3)~Does framing the triangulated value as a
\emph{refinement} of the published initial-value $r_e$ --- rather than as a
result the published derivation must reproduce --- match your reading?}
